% Paragraphs broken with their own tolerances, as in real papers: a
% bibliography (thebibliography sets \sloppy: \tolerance 9999 and
% \emergencystretch 3em) whose entries hold long links and names that do not
% break, a \sloppy paragraph ending in a long link, a paragraph that needs
% \emergencystretch, and one with its own hyphenation penalties.
\documentclass[12pt]{article}
\usepackage{url}
\setlength{\textwidth}{15.5truecm}
\begin{document}
A first paragraph at the default tolerance, long enough to fill a line of the text width and then some more words.

{\sloppy Skorokhod, A.V., Basic principles and applications of probability theory, Springer,2000.\par}

{\sloppy A paragraph set sloppily that ends in a long link. Denial of service via algorithmic complexity attacks. In 12th USENIX Security Symposium (USENIX Security 03). USENIX Association, 2003. URL: \url{https://www.usenix.org/conference/12th-usenix-security-symposium/denial-service-algorithmic-complexity-attacks}.\par}

{\emergencystretch=2em Long words and identifiers such as \texttt{ReflowableDocumentSerializer} and \texttt{paragraph\_tolerance\_parameters} leave a line little to stretch, and the line before them takes the slack.\par}

{\hyphenpenalty=10000 \exhyphenpenalty=10000 \tolerance=1000 A paragraph that does not hyphenate: its well-known, often-cited, state-of-the-art examples of extraordinarily uncomfortable typesetting are characteristically unhyphenated.\par}

\begin{thebibliography}{9}
\bibitem{a}
S.~A. Crosby.
\newblock Denial of service through regular expressions.
\newblock USENIX Association, 2003.
\newblock URL:
  \url{https://www.usenix.org/conference/12th-usenix-security-symposium/denial-service-through-regular-expressions}.

\bibitem{b}
S.~A. Crosby and D.~S. Wallach.
\newblock Denial of service via algorithmic complexity attacks.
\newblock In {\em 12th USENIX Security Symposium (USENIX Security 03)}. USENIX
  Association, 2003.
\newblock \texttt{ReflowableDocumentSerializerConfigurationParameters}.

\bibitem{c}
A.~V. Skorokhod.
\newblock Basic principles and applications of probability theory, Springer,2000.

\bibitem{d}
A last entry.
\end{thebibliography}
\end{document}
